% ID: SEL-D0052-OB-E04
% ESTADO: cotejada. Segunda pasada visual y borrador compilado revisados.
% Fuente: Selectividad/sel_2013_matematicas_aplicadas/septiembre/reserva b/EXAMEN 4-ANDALUCÍA.pdf
% SHA-256: 6201aa5a68f0d1302cfd697751f74c4c73f655f0d9b1ae98ae8960207dbd5225
% Página: 2; opción B; ejercicio original 4.

\begin{enumerate}[label=\alph*)]
\item \textbf{(1.25 puntos)} Se considera la población $\{2,4,6\}$. Escriba todas las posibles muestras de tamaño dos elegidas mediante muestreo aleatorio simple y determine la desviación típica de las medias muestrales.
\item \textbf{(1.25 puntos)} En una ciudad se seleccionó una muestra aleatoria de 500 alumnos de Bachillerato a los que se les preguntó si poseían una determinada marca de teléfono móvil, resultando que 80 de ellos contestaron afirmativamente. Obtenga un intervalo de confianza, al 92\%, para estimar la proporción de estudiantes de Bachillerato que poseen esa marca de teléfono móvil.
\end{enumerate}
